\noindent \textbf{L.4.7 Hadron size: the adopted ring geometry and its limits (v86, P14-C, [T4]).} \emph{Adopted description.} For the charge radii of baryon-type composites the paper adopts the planar Z6 ring as a phenomenological description: psion $k$ at angle $60^\circ k$ on a circle, equal constituent masses, and $\langle r^2\rangle=\sum_kQ_k|\mathbf{x}_k-\mathbf{c}|^2/\sum_kQ_k$ about the centroid (not normalised when $\sum_kQ_k=0$), with one length fixed by the proton ($r_p=0.841$ fm, ring radius $a=1.212$ fm). It is a description, not a derivation from eq.~(1), and it was chosen after a pre-registered series of alternatives (companion repository, WP-GEO to WP-GEO12, WP-HUB, WP-PHS, WP-REL).

\noindent \emph{What it encodes.} For three consecutive labels $(c-1,c,c+1)$ on the ring the mean-square charge radius is the exact identity $\langle r^2\rangle=a^2\big[\tfrac79Q_{\rm tot}-\tfrac23Q_c\big]$. The radii of arc-type triples therefore depend only on cyclic adjacency in Z6 and on the charge of the central label; the star, petal and 12-gon arrangements give identical ratios, and of the 720 orderings of the labels on the ring only the natural cyclic order passes. A neutral triple of mutually non-adjacent labels is an equilateral triangle and has $\langle r^2\rangle=0$ exactly. The neutron radius is negative because the only positive label sits at the centre and the net charge at the rim is negative.

\begin{center}\small
\begin{tabular}{llccl}
State & Triple & Ring $\langle r^2\rangle/\langle r^2\rangle_p$ & Reference & Status \\ \hline
$n$ & $(0,1,5)$ & $-0.154$ & $-0.164$ (measured) & pass (pre-registered) \\
$\Sigma^-$ & $(0,1,2)$ & $+0.846$ & $0.86\pm0.21$ (measured) & see below \\
$\Sigma^0$ & $(0,2,4)$ & $0$ (exact) & $-0.04$ to $+0.10$ (theory) & consistent with zero \\
$\Lambda$ & $(1,2,3)$ & $+0.77$ & $+0.06$ to $+0.15$ (theory) & fails \\
$\Sigma^+$ & $(1,3,5)$ or $(2,3,4)$ & $+2.08$ or $+0.08$ & $+0.84$ to $+1.11$ (theory) & fails \\
exotic & $(3,4,5)$ & $+1.31$ & none observed & open prediction \\
\end{tabular}
\end{center}
\noindent No radius of $\Sigma^+$, $\Sigma^0$ or $\Lambda$ has been measured; the references are chiral-perturbation-theory and lattice results [22,23].

\noindent \emph{Identification caveat.} The $J\!\cdot\!\Sigma$ term couples left to right psions only, so a triple of one chirality has no $J\!\cdot\!\Sigma$ bond. The triple $(0,1,2)$ is then inert, which is the exact statement used for the electron (\S6, L.4.4). The $\Sigma^-$ row treats $(0,1,2)$ as a hadron and is meaningful only if the electron triple and the $\Sigma^-$ are distinct states (different cascade levels, P14-B); it is reported but not counted as an independent success. A model in which the $J\!\cdot\!\Sigma$ bond graph sets the size (a V with the unique-chirality psion as hub, one opening angle fitted to the neutron) makes $(0,1,2)$ pointlike and fails $\Sigma^0$, $\Lambda$ and $\Sigma^+$.

\noindent \emph{Alternatives tested and not adopted.} Line, centred hexagon, blind triangle, prism, spin cone, hourglass, three-triangle star, and a 12-site ring in which $k=3$ is displaced by $30^\circ$ (it fits $\Lambda$ and $\Sigma^+$ by one discrete choice, two of seven positions, with no dynamical origin); the arc and circular-mean metrics of the internal phase circle (neutron $-0.200$). Magnetic moments fail in every point-charge-with-spin rule tested. Neutrino radii cannot be made to vanish together by any placement of the six labels on 12 sites (best $\max|R|=0.056$), so the smallness of the neutrino and of the charged leptons must come from scale, not shape.

\noindent \emph{Limits that remain.} (a) \emph{Mixing.} $J\!\cdot\!\Sigma$ conserves $\sum k$ and mixes $n$, $\Sigma^0$, $\Lambda$ (and $p$ with the two $\Sigma^+$ candidates); the ring assumes no mixing. With adiabatic labels the neutron radius turns positive for $g\gtrsim0.06$ and, taking $p$ and $n$ as the lowest states of their sectors, $m_n<m_p$ at every coupling tested. (b) \emph{Scale.} The number $\hbar c/a\approx163$ MeV is not an energy of the theory: it is $r_p$ converted through the premises that $p=(0,4,5)$, the constituents are classical points of equal mass on one ring radius, and $E=\hbar c/a$ (that is, $\Delta t=a/c$). A massless-fermion confinement correction ($E=x\hbar c/R$, $x=2.04$) raises the energy per psion to 333--349 MeV and gives $M_p=3E$ within $+6\%$ and $+12\%$, but this holds for any three massless constituents in about 1.2 fm, and a common orbital gives $R_n=0$, so the correction can apply only to the length--energy relation. No route available in the paper derives the scale: the $J\!\cdot\!\Sigma$ contact term needs tuning of one part in $10^{24}$; proto-colour running needs $\alpha_{\rm pc}(\Lambda_{\rm meta})\approx1/50$ while every reading in the paper gives $\alpha_{\rm pc}\ge0.09$ (the induced value $3\pi/(S/2g^2)=0.58$ agrees with $\alpha_{\rm meta}=\pi/6$ to 11\%, placing proto-colour confinement near $\Lambda_{\rm meta}$); four pre-registered exponent forms and the cascade rungs miss by factors 3.7 and larger ($\ell=8$ gives $r_p=0.22$ fm, $\ell=8.5$ gives 1.32 fm); a core bound at $\Lambda_{\rm meta}$ is excluded by the proton form factor by 24 orders of magnitude. (c) \emph{Gottfried sum.} The valence part with the charges $Q_k$ is $-1/9$ (quarks $+1/3$), so the sea must supply $+0.346$ to reach the measured $+0.235\pm0.026$, 3.5 times the QCD sea term; a closed list of cloud channels reaches it only for a cloud probability between 0.09 and 0.89, which the PH does not fix, so the proton identification is not excluded.

\noindent This is open problem \textbf{P14-C}: a derived baryon scale, a sea model, and an internal structure consistent with the $J\!\cdot\!\Sigma$ dynamics. The ring is adopted as the working description meanwhile; its failures for $\Lambda$ and $\Sigma^+$ are the first falsifiers, and nothing in the main text depends on this appendix.
